\documentclass[11pt]{article}
\usepackage{geometry}
\usepackage{graphicx}
\usepackage{amsmath}
\usepackage{hyperref}

\geometry{a4paper}

% Commands to easily modify the document for different Stratus projects
\newcommand{\projectname}{XXXXXX}  % Default project name
\newcommand{\caseyear}{XXXX}  % Default year
\newcommand{\caseNumber}{XX} 

% Change this for each experiment
\newcommand{\DeploymentSpike}{2017/05/06 19:00}
\newcommand{\DeploymentRecordedSpike}{2017/05/06 19:05:00}
\newcommand{\RecoverySpike}{2018/04/12 14:50}
\newcommand{\RecoveryRecordedSpike}{2018/04/12 14:45:00}

\title{Data Report for \projectname}
\author{Yu Gao}
\date{\today}

\begin{document}

\maketitle
\tableofcontents
\newpage

\section{Introduction}
This document details the preprocessing, quality control, and analysis steps for the \projectname data, 
capturing significant procedural elements such as Quality Control, data merging, and cataloging.

\section{Pre-Processing}
Detailed steps undertaken during the pre-processing of \projectname\caseNumber data include:


\subsection{Time Marks Check}
The first is obviously that we want the timing to be correct. 
The second is that one of the best ways (if both sensors work for the full deployment) to do quality control is to look at the difference between the two sensors. 
If the time bases are not consistent, there will be a spurious, large variance in the difference.
If the recovery spike changes by more than $\pm$ 5 minutes, appropriate corrections are applied.

% Discussion on the alignment of recorded spike times with observed data
In reviewing the temperature data against recorded spike times for Stratus 16, it is observed
that the differences are consistently within 5 minutes. Specifically, for the deployment phase,
the significant temperature drop commences at \DeploymentSpike, closely aligning with
the recorded spike time of \DeploymentRecordedSpike. 
Similarly, during the recovery phase, a notable temperature drop starts at \RecoverySpike, 
which slightly follows the recorded spike time of \RecoveryRecordedSpike.

% Recommendation based on the findings
Given these findings, interpolation may not be necessary for Stratus 16 as the temperature data 
aligns well with the recorded times, suggesting minimal time offset. Please review these results 
and provide feedback on whether to proceed without interpolation, thus maintaining the original 
data integrity for further analysis.

\begin{figure}[h]
\centering
\includegraphics[width=0.8\textwidth]{/Users/yugao/UOP/ORS-processing/img/\projectname\casenumber_spikes.png}
\caption{Illustration of spike times for \projectname\ \casenumber}
\end{figure}

\begin{figure}[h]
\centering
\includegraphics[width=0.8\textwidth]{/Users/yugao/UOP/ORS-processing/img/\projectname\casenumber_deployment_recovery.png}
\caption{Deployment and recovery phase analysis for \projectname\ \casenumber}
\end{figure}

\subsection{Interpolation}
Data gaps are addressed by referring to the “recover.xls” spreadsheet on buoy5, 
under the “subsurface” tab, to compare spike start and end times with the apparent time in the instrument data. 
If the offset is greater than about 5 minute, data are linearly interpolated to a new time base to synchronize instruments.
No interpolation is needed for \projectname\caseNumber. 


\subsection{Truncation}
Following interpolation, data truncation is performed based on the mooring log. The start time is set to the anchor over time plus 4 hours, and the end time is set to the release fired time.

\subsection{Variables Handling}
Decision on variable inclusion involves whether to carry forward time series for all five variables (temperature, conductivity, salinity, absolute salinity, and pressure) even if not present, by setting missing values to NaN or -999. Variables like pressure are always carried forward.

\section{Quality Control and Deployment Catalog}
The following quality control measures are implemented:

\subsection{Spike Removal}
Spike detection and removal are set up in function "src/qc_function.py". 
The `remove_spikes` function identifies and removes spikes from time series data using a rolling mean and standard deviation. 
It calculates these statistics within a specified window size around each data point. 
Points that deviate from the mean by more than a set multiple of the standard deviation are considered spikes and replaced with NaN. 

\subsection{Difference Statistics}
When both instruments are available and functional, 
the mean and standard deviation of the difference are computed to provide benchmarks for accuracy and precision, 
informing further quality control decisions.

\subsection{Human in the Loop (HITL) Quality Control}
Create a catalog of results for each deployment, featuring original and QC plots for all measured variables. 
These serve as a basis for collaborative review and further data correction.

\section{Merged Data Set}
Procedures for handling and merging the data set:

\subsection{Metadata and Test Runs}
Ensure all necessary metadata are incorporated. Conduct initial test runs to plot data from multiple deployments, 
deciding on integration strategies and overlap handling.

\subsection{Creation of Merged Files}
Generate merged deep temperature data files for each project, detailing the steps in an associated report. 
This includes processing explanations, variable and metadata descriptions, catalog plots, and merged temperature plots.

\section{Documentation and Archiving}
Finalized data sets are documented extensively and prepared for inclusion on the project website. 
Consideration for submission to data repositories with appropriate DOI issuance is also included.

\section{Conclusion}
Summarizes the impact and significance of the preprocessing, quality control, and merging efforts undertaken for the \projectname\caseNumber data.

\end{document}
